\scanpage{25}
\mspage{10}

(3.1), $c_4=0$ si et seulement si $a_2^2=0$. Supposons qu'on ait même $a_2=0$
\[
\begin{aligned}[t]
 &y^2=x^3+a_4 x+a_6 \qquad \pi=a_4^{-1}\,dy=dx/2y\\[.65em]
 &c_4=c_6=0 \qquad \Delta=-a_4^3 \, .
\end{aligned}\tag*{(5.1)}
\]

Les autres systèmes $(\pi,x,y)$ donnant une équation de cette forme s'obtiennent en faisant $s=t=0$ dans (1.2) et
\[
\begin{aligned}
u^4 {\mathbb{\text{\makebox[0pt][l]{\kern.45pt$a$}$a$}}}'_4 &= a_4\\
u^6 a'_6 &= a_6+r\,a_4+r^3 \, .
\end{aligned}\tag*{(5.2)}
\]

b) $p=2$\tikz[baseline=(p.base)]{\node[inner sep=0pt] (p) {.};\handwritten{\draw[thin,gray] (0.13em,1.1em) .. controls (-0.42em,1.25em) and (-0.57em,-1.52em) .. (-0.01em,-1.45em) .. controls (0.40em,-1.46em) and (0.77em,0.33em) .. (0.75em,1.44em);}}

En caractéristique 2 , on a $c_4=a_1^4$. Supposons qu'on ait $a_1=0$ ; $\pi$ étant donnée, on peut alors choisir $x$ et $y$ tels que (1.3) soit
\[
\begin{aligned}[t]
 &y^2+a_3\,y=x^3+a_4\,x+a_6 \qquad
\pi=\frac{dx}{a_3}=dy/\tikz[baseline=(den.base)]{\node[inner sep=0pt] (den) {$x^2+a_4$};\handwritten{\draw[<->,thin] ([yshift=0.6ex]den.north west) .. controls ([xshift=-0.3em,yshift=-1.1ex]den.south west) and ([xshift=-0.1em,yshift=-1.1ex]den.south east) .. ([xshift=0.32em,yshift=-0.35ex]den.south east) .. controls ([xshift=0.6em]den.south east) and ([xshift=0.32em]den.north east) .. ([xshift=0.32em,yshift=0.65ex]den.north east);}}
\\[.65em]
 &b_2=b_4=0 \qquad b_6=a_3^2 \qquad b_8=a_4^2 \qquad
\Delta=a_3^4 \qquad c_4=c_6=0 \, .
\end{aligned}\tag*{(5.3)}
\]

Les autres systèmes $(\tikz[baseline=(pc.base)]{\node[inner sep=0pt] (pc) {$\pi,$};\handwritten{\draw[thin,gray,rotate around={-9:(pc.center)}] (pc.center) ellipse (0.62em and 1.07em);}}\,x,y)$ donnant une équation de cette forme s'obtiennent en faisant $r=s^2$ dans (1.2) et
\[
\begin{aligned}
u^3 a'_3 &= a_3\\
u^4 a'_4 &= a_4+sa_3+s^4\\
u^6 a'_6 &= a_6+s^2\,a_4+ta_3+s^6+t^2 \, .
\end{aligned}\tag*{(5.4)}
\]
% archive sheet number: 24
